
\section{Appendix A. Mathematical Derivations}

\subsection{A.1 k-Means Objective}

Given $N$ points $\{x_i\} \subset \mathbb{R}^d$ and $K$ clusters:

\[
\mathcal{L}(\mathcal{C}, \mu) = \sum_{k=1}^{K} \sum_{i \in \mathcal{C}_k} \|x_i - \mu_k\|_2^2
\tag{A.1}
\]

Lloyd's algorithm monotonically decreases $\mathcal{L}$, converging to a 
local minimum.

\subsection{A.2 Joint Feature Normalization}

We compute:

\[
F_i = \frac{q_i + k_i}{\|q_i + k_i\|_2}
\tag{A.2}
\]

This has unit norm: $\|F_i\|_2 = 1$. k-means on unit-norm features is 
equivalent to spherical k-means.

\subsection{A.3 Cluster-Restricted Attention}

For each cluster $k$, the restricted attention is:

\[
\text{out}_i^{(k)} = \begin{cases} 
\text{softmax}_{j \in \mathcal{S}_k}\!\left(\frac{q_i^\top k_j}{\sqrt{d}}\right) V_{\mathcal{S}_k} & i \in \mathcal{S}_k \\
0 & \text{otherwise}
\end{cases}
\tag{A.3}
\]

where $\mathcal{S}_k = \{j : c_j = k\} \cup \{0\}$.

\subsection{A.4 Complexity of Naive Implementation}

The loop over $K$ clusters, each computing a full $L \times L$ matrix:

\[
T_\text{naive} = K \cdot \left( L^2 d + L^2 + L^2 d \right) = \mathcal{O}(K L^2 d)
\tag{A.4}
\]

Compare to dense: $T_\text{dense} = \mathcal{O}(L^2 d)$. Ratio: $K$.

\textbf{Theoretical sparse cost} (with sparse matmul):

\[
T_\text{sparse} = \mathcal{O}(L \cdot (L/K) \cdot d) = \mathcal{O}(L^2 d / K)
\tag{A.5}
\]

Savings: $K$ if sparse ops are used.

\subsection{A.5 Non-Differentiability of argmin}

The \texttt{argmin} operation is piecewise constant: it has zero gradient almost 
everywhere and undefined gradient at boundaries. Thus:

\[
\frac{\partial c_i}{\partial q_i} = 0 \text{ a.e.}
\tag{A.6}
\]

No gradient flows through cluster assignment.

\textbf{Alternative.} Gumbel-Softmax relaxation:

\[
c_i^{(k)} = \frac{\exp(-\|f_i - \mu_k\|_2^2 / \tau)}{\sum_{k'} \exp(-\|f_i - \mu_{k'}\|_2^2 / \tau)}
\tag{A.7}
\]

which is differentiable for all $\tau > 0$.

\subsection{A.6 Token-0 Anchor}

Token 0 is included in every cluster. Its effect is to provide a shared 
context key. In terms of attention weights:

\[
\text{attn}_{i,0} = \frac{e^{q_i^\top k_0 / \sqrt{d}}}{\sum_{j \in \mathcal{S}_{c_i}} e^{q_i^\top k_j / \sqrt{d}}}
\tag{A.8}
\]

If token 0 is a "sink" (Xiao et al., 2024), $\text{attn}_{i,0}$ is 
consistently large, providing a stable baseline for the softmax denominator.
\begin{lstlisting}

---
\end{lstlisting}